% Ringville, and the two other figures the Module 2 sheet draws. Kept in their
% own file because the answer key draws exactly the same town with the boxes
% filled in. Needs tikz, its calc and positioning libraries, and the colours
% pltblue and pltrust, all of which the including file sets up.
% ---------------------------------------------------------------------------
% The town: sixteen people in a circle, everybody friends with the two on
% their left and the two on their right. Four friends each, thirty-two
% friendships. That one drawing has to carry both halves of the module --
% distance AND triangles -- which is why it is next-two-neighbours rather than
% next-one: a ring of next-one neighbours has no triangle anywhere in it.
%
% Two more things fall out of the choice and both are used on the sheet. The
% town is left-right symmetric about the line through person 1 and person 9,
% so a student who has worked out people 2 to 9 already has 10 to 16. And the
% two shortcuts (1--9 and 5--13) are diameters, so they cross at the middle of
% the drawing and read as two lines rather than as a tangle.
% ---------------------------------------------------------------------------
\tikzset{
  % Every friendship is drawn the same, so no line weight says "this one
  % matters more". The two shortcuts are the exception, and they are the point.
  friend/.style   = {line width=1.4pt, black!58, line cap=round},
  shortcut/.style = {line width=2.1pt, pltrust, line cap=round},
  chair/.style    = {circle, draw=black!72, fill=white, line width=0.9pt,
                     minimum size=0.66cm, inner sep=0pt, font=\footnotesize},
  chairhot/.style = {chair, draw=pltblue!75!black, line width=1.5pt,
                     fill=pltblue!14},
  % The box a student writes a handshake count into, one per person, printed
  % outside the ring where no friendship line can reach it.
  dbox/.style     = {draw=pltblue!70, line width=0.7pt, fill=white,
                     minimum size=0.58cm, inner sep=0pt, rounded corners=1.2pt},
  dnum/.style     = {font=\bfseries\footnotesize, inner sep=0pt},
  % A pair that might or might not be a friendship. The student darkens it.
  maybe/.style    = {line width=1.3pt, black!52, dash pattern=on 1.6pt off 2.2pt},
}

% \ringtown{extra lines}{numbers already written in}.
%
% The count goes in the box beside the person, on the drawing, rather than
% into a list of blanks underneath it: the thing the sheet is teaching is that
% the number belongs to a seat, and a column of fifteen blanks hides that.
\newcommand{\ringtown}[2]{%
\begin{tikzpicture}[x=1cm, y=1cm]
    \foreach \i in {1,...,16} {
        \pgfmathsetmacro{\ang}{90 - (\i-1)*22.5}
        \coordinate (c\i) at (\ang:2.7);
    }
    % Next door, both sides: the rim of the circle.
    \foreach \i in {1,...,16} {
        \pgfmathtruncatemacro{\j}{mod(\i,16)+1}
        \draw[friend] (c\i) -- (c\j);
    }
    % One seat over, both sides. Bowed inward: drawn straight, the chord passes
    % 0.14cm inside the chair sitting between its two ends and is drawn through
    % that person's face.
    \foreach \i in {1,...,16} {
        \pgfmathtruncatemacro{\j}{mod(\i+1,16)+1}
        \pgfmathsetmacro{\amid}{90 - \i*22.5}
        \draw[friend] (c\i) .. controls (\amid:2.12) .. (c\j);
    }
    #1
    \foreach \i in {1,...,16} {
        \pgfmathsetmacro{\ang}{90 - (\i-1)*22.5}
        \node[chair] (p\i) at (\ang:2.7) {\i};
        \node[dbox, name=b\i] at (\ang:3.44) {};
    }
    #2
\end{tikzpicture}}

% The two friendships that came back from college. Both are diameters of the
% circle -- 1 is opposite 9 and 5 is opposite 13 -- so they cross at the
% centre, which is the one place in the drawing with nothing else in it.
\newcommand{\shortcuts}{%
    \draw[shortcut] (c1) -- (c9);
    \draw[shortcut] (c5) -- (c13);}

% A number written into person #1's box.
\newcommand{\dfill}[2]{\node[dnum] at (b#1) {#2};}

% Every distance from person 1 in the plain town, and in the town with the two
% shortcuts. The sheet prints two of each; the answer key prints them all.
\newcommand{\plainfill}{%
  \dfill{1}{0}\dfill{2}{1}\dfill{3}{1}\dfill{4}{2}\dfill{5}{2}\dfill{6}{3}%
  \dfill{7}{3}\dfill{8}{4}\dfill{9}{4}\dfill{10}{4}\dfill{11}{3}\dfill{12}{3}%
  \dfill{13}{2}\dfill{14}{2}\dfill{15}{1}\dfill{16}{1}}
\newcommand{\cutfill}{%
  \dfill{1}{0}\dfill{2}{1}\dfill{3}{1}\dfill{4}{2}\dfill{5}{2}\dfill{6}{3}%
  \dfill{7}{2}\dfill{8}{2}\dfill{9}{1}\dfill{10}{2}\dfill{11}{2}\dfill{12}{3}%
  \dfill{13}{2}\dfill{14}{2}\dfill{15}{1}\dfill{16}{1}}

% ---------------------------------------------------------------------------
% Person 3's four friends, lifted out of the circle. Person 3 is NOT in this
% picture: drawn in, their four lines cross the six pairs the student is being
% asked to read, and the figure stops being readable at exactly the moment it
% is needed.
%
% The four are placed 1, 4, 2, 5 clockwise from the top left rather than in
% seat order. In seat order the three real friendships are the three shortest
% lines in the figure and the answer can be had without looking at the town;
% shuffled, the student has to go back to the circle for all six.
% ---------------------------------------------------------------------------
\newcommand{\hoodfig}[1]{%
\begin{tikzpicture}[x=1cm, y=1cm]
    \coordinate (h1) at (-1.55, 1.30);
    \coordinate (h4) at ( 1.55, 1.30);
    \coordinate (h2) at ( 1.55,-1.30);
    \coordinate (h5) at (-1.55,-1.30);
    \draw[maybe] (h1) -- (h4);
    \draw[maybe] (h4) -- (h2);
    \draw[maybe] (h2) -- (h5);
    \draw[maybe] (h5) -- (h1);
    \draw[maybe] (h1) -- (h2);
    \draw[maybe] (h4) -- (h5);
    #1
    \foreach \i in {1,2,4,5} {\node[chair] at (h\i) {\i};}
\end{tikzpicture}}

% The three of the six that really are friendships, for the answer key.
\newcommand{\hoodreal}{%
    \draw[friend, pltblue!80!black] (h4) -- (h2);
    \draw[friend, pltblue!80!black] (h1) -- (h2);
    \draw[friend, pltblue!80!black] (h4) -- (h5);}
